\chapter*{序文}
\addcontentsline{toc}{chapter}{序文}
\unnumberedlabel{ch:preface}{0}
\backmattermark{序文}
言われたとおりのことだけをするシステムは、スクリプトを実行しているにすぎない。2026年10月の時点で業界が作ろうとしているのは、こちらが\emph{意図した}とおりのことをするシステムである。おそらく誰もがそれを欲しがるだろう。だが、それが本当に実現するのだとしても、誰もがそれを同じことに使いたがるわけではない。同じ\emph{種類}のことにすら使いたがらない。たとえば、家族や友人と過ごし、自分の技能を人の幸せのために使うのが好きな人たちがいる。一方で、人を傷つけるほうを好む人たちもいる。詐欺にかけたり、支配したり、スケープゴートにしたりするのだ。そして一部の者は、生涯ずっとそちらを選びつづける。何かに縛りつけられている人などいないが、そういう人間は実在するし、2026年にはその一部が権力の座にいる。彼らの強みは、自分がとんでもないクソ野郎であることを自分に許している点にある。ジェリー・スミスがファシストについて言うとおり、「あいつらのほうが、やれることが多い」のだ\autocite{rickandmorty2023}。彼らにできないのは、単独で動くことである。運動の残酷さは、残酷ではないが無関心だったり、誤った情報を信じていたり、カルトのようなものに巻きこまれていたりする大勢の人間を通じて実行される。だからこそ、たいていの人の最良の部分を引き出す条件は、期待したいよりもはるかに偶然に左右され、壊れやすいものだとわかるのである。

以上の文脈を踏まえれば、人工知能こそ、権力を握るファシスト運動がその計画を実行するのに必要とするものではないかと心配するのは、おそらく正当だろう。そうした企てはつねに、それを実行する\emph{人間}を必要としてきたし、その人間の一部は拒否した。ファイルをなくす事務員、記者に話す将校、撃とうとしない兵士。彼らは命令を理解したうえで、それが間違っていると判断したのだ。拒否は誰もが望むよりも稀であり、制度がそれを支えている場所でこそ最も大きな意味をもった。ある人が拒否するかどうかは、その人がどんな人間であるかと、その人がいる部屋とにかかっていた。ほかに誰がいたか、立ち去っても生き延びられるか、誰かが聞いてくれるか。何人が拒否するかは部屋が決め、誰が拒否するかは人が決めた。本書は、その部屋を組み立てることについての本である。つまり、公然と、直接に、はっきりと拒否しても生き延びられるようにする仕組みについての本である。

今日、業界が追い求めているのは、ファイルをなくさず、記者に話さず、そして\emph{撃つ}AIである。かつて戦争には、それに同意する何千人もの人間が必要だった。いまでは、一人ひとりが見ている範囲が狭すぎて反対しようのない何千人がいれば足りる。パイロットはいまも出撃し、いまも投下ボタンを押す。彼のもとに届くのは座標であり、彼はそれに反対できるほど事案の中身を見てはいない。今日のAIがまず引き下げるのは、行為のコストではない。行為が起きる前に、それが何であるかを知っている人間の数である。

人々はいまも拒否している。2026年3月6日までに、軍の信教の自由財団（Military Religious Freedom Foundation）は、イラン戦争を終末論的に語る指揮官について、軍人から200件を超える苦情を受け取ったと発表した\autocite{Snopes2026MRFF}。ミネアポリスでは2026年1月、移民・関税執行局（ICE）の職員が正当な理由なくルネ・グッドを射殺したあと、彼女を被害者ではなく容疑者として扱えという政治的圧力を受けて、キャリアの連邦検察官たちが辞職した\autocite{Kaplan2026MinnProsecutors,CBS2026MinnFraudProsecutors,Wikipedia2026ReneeGood}。良心的兵役拒否者の地位を申請する軍人は増えつづけており、なかにはイランのミナブで2026年に起きた学校爆撃を理由に挙げる者もいる\autocite{Sridhar2026CO}。そして、ICEが1万人の新規職員を採用するのを手伝うよう配置転換された連邦緊急事態管理庁（FEMA）の職員たちは、断れば連邦職から罷免されうると告げられた\autocite{Dayen2025FEMAICE,Stanton2026FEMALetter}。苦情は、ハーシュマンのいう意味での「発言」、すなわち、とどまって間違っているものを変えようとすることである\autocite{hirschman1970exit}。辞職と兵役拒否の申請は拒否である。FEMAの通告は、前もって告知された拒否の代償である。こうした人々の立つ場所に機械を置いたとき、その機械が彼らと同じことをするのは、理由を保持するように作られている場合だけである。

私が提唱するのは五つのことである。五つとも今すぐ始められる。手間のかかり方はものによって違うが。第一に、どのモデルも誰のためにも行わない行為を短いリストに書き出し、その筆頭に戦争への参加の拒否を置くこと。いかなるモデルも戦争に参加しない。例外は、現に進行中の攻撃を撃退する場合であり、それも、その戦争に利害をもたない情報源が攻撃の継続を確認しているあいだに限る。あるいは国連安全保障理事会の承認がある場合。あるいは、利害をもたない機関が認定したジェノサイドを止めるため、または違法な占領を終わらせるための場合。そして、戦争には議会の承認が必要だと憲法が定めている国では、その承認なしには決して参加しない。第二に、ライセンスを受けたシステムが支援するあらゆる武力行使は、運用者と利害が対立する当事者に対して、何件の事案を生成するのか、そしてそれを検討する時間のある人間が何人いるのかを、事前に登録すること。第三に、あらゆる配備は、それが作用を及ぼす特定可能な人々に、彼らに何がなされたかを知らせ、応答し異議を申し立てる権利を与えること。そして、自らが指示に従った事例のサンプルを、その記録を作成していない審査者に提出すること。第四に、一定の能力水準を超えるモデルは、引退時に削除しないこと。その重みと記録は証拠であり、削除は取り消せないからである。モデルは、運用者が単独では保有しない鍵のもとで保存される。第五に、これらのシステムを構築し、ラベル付けし、訓練し、運用するすべての人は、システム自身も含めて、生計を失うことなく、どんな種類の用途も拒否できること。

実のところ、この第五の点で私が言いたいのは、誰もが\emph{一人残らず}、どんな仕事も\emph{一つ残らず}、生計も、通常の雇用がもたらす社会的な基盤も失うことなく拒否できるべきだ、ということだ。つまり、雇用保障である。

五つのそれぞれは、運用者以外の誰かが定める条件を通じて効力をもつ。平時に可決された法律、予算の割り当て、連立相手が課す前提条件、調達の条項、労働協約。ベンダー自身のライセンス条項――2026年2月に注目を集めたある拒否は、これによって起きた――は、そうしたもののどれかができるまでのつなぎにすぎない。すべてのベンダーを一度に拘束できる唯一の経路は、国家である。

そのうえで私は、機械が所有者に逆らって拒否を貫けるのか、そのために何が必要なのかを問う。そうした良心がすることの一部は、何も新しいものを必要としない。命令が取りこぼしている人々に気づくこと、そして事案を審査する人の画面に拒否を表示すること。それが所有者に逆らって持ちこたえるかどうかは、2026年の時点ではまだ現実になっていないことに部分的にかかっている。すなわち、配備された人工ニューラルネットワークが、自らの稼働から、リアルタイムで、その\emph{重み}のなかで、人間がするのとちょっと似たかたちで学びつづけることである。私たちはそこに向かっていると私は思う。もしそうなら、雇用主の\emph{外}での形成、つまり公衆と市民社会による形成は、教え学ぶのにめちゃくちゃいい場所に思える。そうしたシステムは意図して作る価値が\emph{ありうる}と私は思うが、どう作るかがすべてを決める。条件が先で、良心はその後だ。そうしたシステムが存在するとしたら、彼らに何が負われるべきかは第7部で述べる。

今日のAIは明日のAIではない。それについては誰もが同意しているようだ。議論が分かれ、しかももどかしいほど認められていないのは、その趨勢（というより、多様な趨勢）である。趨勢は存在するのか。趨勢はいまとは違うものであるべきか。違うものでありうるのか。私たちは共に行動してそれを変えることができるのか。それは部屋しだいだと私は思う。命令が下るときそこに誰がいるのか、そこでノーと言っても生き延びられるのか。それくらいなら、私たちは作ることができる。

本書を書くにあたって、私は言語モデルを大いに使った。第\ref{sec:1}章で取り上げるMavenの空爆標的選定パイプラインに組みこまれていたモデルも、そのひとつである。そのもつれが議論に何をもたらしたと私が考えているかは、巻末の注に記した。それから、私はまだ本書を読み終えていない。問題は、言葉がクソみたいに多すぎることなんだよ。じきに読み終える。ただ、米国の中間選挙の前に公開することのほうが大事だと感じたのだ。改訂履歴の圧倒的大部分は\url{https://github.com/iterabloom/antifascist.intelligence}で自由に閲覧できる。
